\begin{textsample}{POS Dim 3 – global\_south\_2024\_08 – Score 18.00 – global\_south\_2024\_08/112445675.txt}  \label{ex:f3_pos_301}
Returning Gaza \textbf{Student} \textbf{Protesters} Face Punitive Welcome Share : As \textbf{university} and \textbf{college} \textbf{students} prepare to return to \textbf{campus} for the new \textbf{school} year, many find themselves in legal limbo for their participation in the anti-Genocide \textbf{student} \textbf{encampments} \textbf{protesting} Israel ' s war in Gaza at the end of the last \textbf{school} year.
On August 14, the Council on American-Islamic Relations ( CAIR ), Arab American Anti-Discrimination Committee ( ADC ) and Palestine Legal held a press conference to discuss how \textbf{universities} are punishing returning \textbf{students} for their participation in the \textbf{protests}.
CAIR today named three \textbf{universities} as"\textbf{institutions} of particular concern,"ie.
George Washington \textbf{University} ( GW ) in Washington, DC, \textbf{University} of California, Los Angeles ( UCLA ) and Emory \textbf{University} in Atlanta, Georgia.
As an example of punishment for their involvement in \textbf{protests}, some \textbf{student} \textbf{protesters} who were \textbf{arrested} at GW, according a Washington Post article on August 10, would have to agree to accept certain conditions in order to return and"... may be able to have their charges dropped -- but only @ @ @ @ @ @ @ @ @ @ to \textbf{campus} for six months."The list of strict conditions was reported as"not being able to enter dining halls, nor study sessions in the library.
No meeting up with friends for coffee."And that,"\textbf{students} would only be allowed to go to and from their \textbf{residence} and \textbf{classes}, with exceptions for accessing the hospital or using the metro."The conditions as set forth in the agreement constitute a form of"house \textbf{arrest}"while using the term"may"for the dropping of charges, still leaves \textbf{students} with no guarantee that will in fact be the case.
Additionally, the condition of denying such access to educational resources for some \textbf{students} at the private \textbf{university} who pay tuition as high as \$70,000 a year is both punitive and discriminatory in nature and without merit.
In disallowing access to those resources, the \textbf{university} has created two \textbf{classes} of \textbf{students} by discriminating against those who have participated in the \textbf{protests}.
Dylan Sabah, staff \textbf{attorney} for Palestine Legal @ @ @ @ @ @ @ @ @ @ pressure from donors and outside lobby groups have sought to suppress Palestine advocacy on their \textbf{campuses}."And in so doing, he iterated how civil and constructional rights of \textbf{students} have been violated.
Going further, he told of how some \textbf{universities} have allowed"outside actors"to gain access to their \textbf{campuses} to harass and"dox"\textbf{students} and that some have banned entire \textbf{student} groups including Jewish and Palestinian as well as others for engaging in"protected political expression in demanding an end to the genocide and an end to support for Israel."Out of \textbf{student} frustration of not being heard by administrations, \textbf{students} began their \textbf{protests} in setting up \textbf{encampments} to demand a ceasefire and an end to the \textbf{university} ' s complicity in the genocide.
Those actions in turn saw administrations respond by calling in \textbf{police} to break up the \textbf{encampments}, with \textbf{police} brutalizing and \textbf{arresting} \textbf{students}, \textbf{faculty} and staff.
Sabah commented further noting that that reaction is a"profound embarrassment to these \textbf{universities} who pretend to support free expression @ @ @ @ @ @ @ @ @ @ said the \textbf{universities}"... remain committed in silencing dissent and protecting the interest of the powerful against the righteous anger of the many."He suggested that \textbf{institutions} should join in with their \textbf{students} to be"on the right side of history."Chris Godshall-Bennett, Legal Director of the Arab American Anti-Discrimination Committee, accused \textbf{universities} of a"disturbing phenomenon of \textbf{students} begin leveled with extreme discipline charges... prior to any \textbf{hearing} being conducted."He named GW as a"particularly brutal example"of punishment where \textbf{students} were brought before \textbf{student} conduct panels having already had their housing \textbf{revoked} in advance.
He followed by describing those \textbf{students} engaged in the \textbf{protests} and \textbf{encampments} as the"Beating heart of American solidarity for Palestine."In closing, Corey Saylor, CAIR Research and Advocacy Director, announced that its"Unhostile \textbf{Campus} Campaign,"titled"Hostile"concerning the"Targeting of Anti-Genocide \textbf{protesters} while enabling Anti-Palestinian Racism and Islamophobia"has been mailed to 600 \textbf{university} \textbf{administrators} across the country. @ @ @ @ @ @ @ @ @ @ role that \textbf{universities} have played in targeting \textbf{students} \textbf{protesting} the Israeli government ' s ongoing genocide against Palestinians in Gaza.
He suggested that"While \textbf{campuses} welcome back \textbf{students} we should celebrate and protect such \textbf{students}, not focus on harsh punishments for \textbf{protesters}."Although today ' s announcement covered only three \textbf{universities}, Saylor promised that in the weeks ahead CAIR would be adding other \textbf{institutions} of higher learning to their"\textbf{institutions} of particular concern"list.

% matched lemmas: administrator, arrest, attorney, campus, class, college, encampment, faculty, hearing, institution, police, protest, protester, residence, revoke, school, student, university
\end{textsample}
